% Options for packages loaded elsewhere
\PassOptionsToPackage{unicode}{hyperref}
\PassOptionsToPackage{hyphens}{url}
%
\documentclass[12pt]{article}
\usepackage[left=2cm, right=2cm, top=2cm, bottom=2cm]{geometry}
\usepackage{amsmath,amssymb}
\usepackage{iftex}
\ifPDFTeX
  \usepackage[T1]{fontenc}
  \usepackage[utf8]{inputenc}
  \usepackage{textcomp} % provide euro and other symbols
\else % if luatex or xetex
  \usepackage{unicode-math} % this also loads fontspec
  \defaultfontfeatures{Scale=MatchLowercase}
  \defaultfontfeatures[\rmfamily]{Ligatures=TeX,Scale=1}
\fi
\usepackage{lmodern}
\ifPDFTeX\else
  % xetex/luatex font selection
\fi
% Use upquote if available, for straight quotes in verbatim environments
\IfFileExists{upquote.sty}{\usepackage{upquote}}{}
\IfFileExists{microtype.sty}{% use microtype if available
  \usepackage[]{microtype}
  \UseMicrotypeSet[protrusion]{basicmath} % disable protrusion for tt fonts
}{}
\makeatletter
\@ifundefined{KOMAClassName}{% if non-KOMA class
  \IfFileExists{parskip.sty}{%
    \usepackage{parskip}
  }{% else
    \setlength{\parindent}{0pt}
    \setlength{\parskip}{6pt plus 2pt minus 1pt}}
}{% if KOMA class
  \KOMAoptions{parskip=half}}
\makeatother
\usepackage{xcolor}
\usepackage{longtable,booktabs,array}
\usepackage{calc} % for calculating minipage widths
% Correct order of tables after \paragraph or \subparagraph
\usepackage{etoolbox}
\makeatletter
\patchcmd\longtable{\par}{\if@noskipsec\mbox{}\fi\par}{}{}
\makeatother
% Allow footnotes in longtable head/foot
\IfFileExists{footnotehyper.sty}{\usepackage{footnotehyper}}{\usepackage{footnote}}
\makesavenoteenv{longtable}
\usepackage{graphicx}
\makeatletter
\def\maxwidth{\ifdim\Gin@nat@width>\linewidth\linewidth\else\Gin@nat@width\fi}
\def\maxheight{\ifdim\Gin@nat@height>\textheight\textheight\else\Gin@nat@height\fi}
\makeatother
% Scale images if necessary, so that they will not overflow the page
% margins by default, and it is still possible to overwrite the defaults
% using explicit options in \includegraphics[width, height, ...]{}
\setkeys{Gin}{width=\maxwidth,height=\maxheight,keepaspectratio}
% Set default figure placement to htbp
\makeatletter
\def\fps@figure{htbp}
\makeatother
\setlength{\emergencystretch}{3em} % prevent overfull lines
\providecommand{\tightlist}{%
  \setlength{\itemsep}{0pt}\setlength{\parskip}{0pt}}
\setcounter{secnumdepth}{-\maxdimen} % remove section numbering
\ifLuaTeX
  \usepackage{selnolig}  % disable illegal ligatures
\fi
\IfFileExists{bookmark.sty}{\usepackage{bookmark}}{\usepackage{hyperref}}
\IfFileExists{xurl.sty}{\usepackage{xurl}}{} % add URL line breaks if available
\urlstyle{same}
\hypersetup{
  hidelinks,
  pdfcreator={LaTeX via pandoc}}

\author{}
\date{}

\begin{document}

\protect\hypertarget{page-0-0}{}{}EI SEVIER


\hypertarget{applied-mathematics-and-computation}{%
\subsubsection{\texorpdfstring{\textbf{Applied Mathematics and
Computation}}{Applied Mathematics and Computation}}\label{applied-mathematics-and-computation}}

\hypertarget{numerical-solution-of-one-dimensional-sinegordon-equation-using-high-accuracy-multiquadric-quasi-interpolation}{%
\subsection{Numerical solution of one-dimensional Sine--Gordon equation using high-accuracy multiquadric quasi-interpolation
★}\label{numerical-solution-of-one-dimensional-sinegordon-equation-using-high-accuracy-multiquadric-quasi-interpolation}}

Zi-Wu Jiang *, Ren-Hong Wang

School of Mathematical Sciences, Dalian University of Technology, Dalian
116024, PR China \\

E-mail address: zwjiang@yahoo.cn (Z.-W. Jiang).

† This work was supported by the National Natural Science Foundation of China (U0935004, 110 71031, 11001037 and 10801024).

* Corresponding author.

\hypertarget{abstract}{%
\paragraph{ABSTRACT}\label{abstract}}

In this paper, we propose a numerical scheme to solve the one-dimensional Sine-Gordon equation related to many scientific research topics by using high-accuracy multiquadric quasi-interpolation. We use the derivatives of a multiquadric quasi-interpolant to approximate the spatial derivatives, and a finite difference to approximate the temporal derivative. The advantages of the scheme are that it is meshfree, and in each time step only a multiquadric quasi-interpolant is employed, so that the algorithm is very easy to implement. The accuracy of our scheme is demonstrated by some test problems.

© 2012 Elsevier Inc.~All rights reserved.

\hypertarget{introduction}{%
\paragraph{1. Introduction}\label{introduction}}

The nonlinear one-dimensional Sine-Gordon (SG) equation

\[\frac{\partial^2 u}{\partial t^2} = \frac{\partial^2 u}{\partial x^2} - \sin(u), \quad a \leqslant x \leqslant b, \ t \geqslant 0, \tag{1}\]

which appear in differential geometry gained its significance because of the collisional behaviors of solitons that arise from these equations.
This equation appeared in many scientific fields such as the propagation of fluxion in Josephson junctions {[}15{]} between two superconductors, the motion of rigid pendular attached to a stretched wire, solid state physics, nonlinear optics, and dislocations in metals {[}19{]}. This paper is devoted to the numerical solution of equation (1) with the initial conditions

\[u(x,0) = f(x), \quad a \le x \le b,\tag{2}\] 

\[\frac{\partial u}{\partial t}(x,0) = g(x), \quad a \leqslant x \leqslant b,\tag{3}\]

and the boundary conditions

\[u(a,t) = h_1(t), \quad u(b,t) = h_2(t), \quad t \geqslant 0.\tag{4}\]

Since the exact solution of the SG equation can only be obtained in special situations, many numerical schemes are constructed to solve the SG equation. To mention some of them: finite difference schemes {[}3,4,8,11{]}, finite elements method {[}1{]}, and B-spline collocation method {[}16{]} etc. Although the previous methods have been widely used, they usually require the construction and update of a mesh and hence bring inconvenience during computation. To avoid the mesh generation, in recent years, meshless techniques have attracted attention of researchers. Dehghan and Shokri {[}7{]} proposed a numerical scheme to solve one-dimensional SG equation using collocation points and approximating directly the solution using the thin plate splines radial basis function (RBF). Unfortunately, this scheme need to solve a ill-conditional system of linear equations. Ma and Wu {[}13{]} presented a meshless scheme by using a multiquadric (MQ) quasi-interpolation named \(\mathcal{L}_{\mathcal{E}}\) without solving a large-scale linear system of equations, but a polynomial p(x) was needed to improved the accuracy of the scheme. In this paper we developed a meshless approach by directly using high accuracy MQ quasi-interpolation defined on {[}a,b{]} without using any polynomial.

The remainder of this paper is organized as follows: In Section 2, some elementary knowledge about the RBF interpolation and the MQ quasi-interpolation are introduced. In Section 3, we present the numerical techniques by using the high accuracy MQ quasi-interpolation to solve the SG equation. The computational results of the SG equation for some test problems are illustrated, and compared with those obtained with previous results in Section 4.

\hypertarget{rbf-interpolation-and-mq-quasi-interpolation}{%
\paragraph{2. RBF interpolation and MQ quasi-interpolation}\label{rbf-interpolation-and-mq-quasi-interpolation}}

\hypertarget{rbf-interpolation}{%
\paragraph{2.1. RBF interpolation}\label{rbf-interpolation}}

RBF was introduced by Krige {[}10{]} in 1951 to deal with geological problem. Hitherto, RBFs were widely used in many fields ({[}6,18{]}, etc). The process of interpolation by using RBF is as follows.

For a given region \(\Omega \in \mathbb{R}^d\) and a set of finite distinct interpolation points

\[X = \{x_1, \ldots, x_N\} \subset \Omega\]

if we are supplied with a function \(f: X \to \mathbb{R}\) , we can construct an interpolant to f in the form:

\[S_{f,X}(x) = \sum_{i=1}^{N} \alpha_j \varphi(\|x - x_j\|), \quad \text{for } x \in \Omega,\tag{5}\]

where \(\|\cdot\|\) denote Euclidean norm, and \(\varphi\) is a certain RBF, say \(\varphi:\mathbb{R}_{\geqslant 0}\to\mathbb{R}\) . The undetermined coefficients \(\{\alpha_1,\ldots,\alpha_N\}\) can be obtained by the linear system

\[\sum_{i=1}^{N} \alpha_j \varphi(\|\mathbf{x}_i - \mathbf{x}_j\|) = f(\mathbf{x}_i), \quad 1 \leqslant i \leqslant N.\tag{6}\]

In this paper, inverse multiquadric (IMQ) RBF is chosen as the RBF mentioned in (5), and the adopted IMQ-RBF has the form:

\[\varphi(r) = \frac{s^2}{(s^2 + r^2)^{3/2}},\tag{7}\]

where \emph{s} is a shape parameter. From {[}14{]}, we know that the linear system (6) have unique solution.

\hypertarget{mq-quasi-interpolation}{%
\paragraph{2.2. MQ quasi-interpolation}\label{mq-quasi-interpolation}}

Quasi-interpolation of a univariate function \(f:[a,b]\to\mathbb{R}\)
with MQs on the scattered points

\[a = x_0 < x_1 < \dots < x_n = b, \quad h := \max_{1 \le i \le n} (x_i - x_{i-1}),\tag{8}\]

usually takes the form

\[(\mathcal{M}f)(x) = \sum_{i=0}^{n} f(x_i)\psi_i(x), \tag{9}\]

where \(\psi_i(x)\) is linear combinations of the MQs

\[\phi_i(x) = \sqrt{c^2 + (x - x_i)^2}  \tag{10}\] 

 \[ c \in \mathbb{R}^+ \]is a shape parameter.

In 1988, Buhmann {[}5{]} studied the accuracy of quasi-interpolation operator \(\mathcal{M}\) of a function
\(f: \mathbb{R} \mapsto \mathbb{R}\) at infinite regular points in \(\mathbb{R}\) and showed that this scheme reproduces linear
polynomials. Beatson and Powell {[}2{]} proposed three univariate MQ quasi-interpolations, namely, \(\mathcal{L}_{\mathcal{A}}\) , \(\mathcal{L}_{\mathcal{B}}\) , and \(\mathcal{L}_{\mathcal{C}}\) on interval {[}a,b{]}. Wu and Schaback {[}20{]} proposed the univariate quasi-interpolation \(\mathcal{L}_{\mathcal{D}}\) , and proved that the scheme possess preserving monotonicity and linear reproducing property as well as \(\mathcal{O}(h^2)\) error of quasi-interpolant
\(\mathcal{L}_{\mathcal{D}}f\) in the case of
\(c^2|\log c|=\mathcal{O}(h^2)\) . Based on Wu and Schaback's work, Ling {[}12{]} proposed a multilevel univariate quasi-interpolation \(\mathcal{L}_{\mathcal{R}}\) , and this scheme can converge with a rate of \(\mathcal{O}(h^{2.5}\log h)\) as \(c=\mathbb{O}(h)\) . To solve SG equation numerically, in this paper we adopt a MQ quasi-interpolation scheme constructed in {[}9{]} using multilevel method, named \(\mathcal{L}_{\mathcal{W}_2}\) .

Next we recall the construction of \(\mathcal{L}_{w_2}\) based on the IMQ-RBF interpolation and the \(\mathcal{L}_{\mathcal{D}}\) operator.
Given data \(\{(x_i, f_i)\}_{i=0}^n\) , Wu--Schaback's
\(\mathcal{L}_{\mathcal{D}}\) operator is

\[(\mathcal{L}_{\mathcal{D}})f(x) = \sum_{j=0}^{n} f_j \psi_j(x), \tag{11}\]

where

\[\psi_{0}(x) = \frac{1}{2} + \frac{\phi_{1}(x) - (x - x_{0})}{2(x_{1} - x_{0})}, \\
\psi_{1}(x) = \frac{\phi_{2}(x) - \phi_{1}(x)}{2(x_{2} - x_{1})} - \frac{\phi_{1}(x) - (x - x_{0})}{2(x_{1} - x_{0})}, \\
\psi_{n-1}(x) = \frac{(x_{n} - x) - \phi_{n-1}(x)}{2(x_{n} - x_{n-1})} - \frac{\phi_{n-1}(x) - \phi_{n-2}(x)}{2(x_{n-1} - x_{n-2})}, \\
\psi_{n}(x) = \frac{1}{2} + \frac{\phi_{n-1}(x) - (x_{n} - x)}{2(x_{n} - x_{n-1})}, \\
\psi_{j}(x) = \frac{\phi_{j+1}(x) - \phi_{j}(x)}{2(x_{j+1} - x_{j})} - \frac{\phi_{j}(x) - \phi_{j-1}(x)}{2(x_{j} - x_{j-1})}, \quad j = 2, \dots, n - 2.\tag{12}\]

Suppose that \(\{x_i\}_{i=0}^n\) are given scattered points, and N is a positive integer satisfying N \textless{} n, then we define an index \(0 < k_1 < k_2 < \cdots < k_N < n\) , and pick a smaller set denoted by
\(\{x_{k_j}\}_{j=1}^N\) from the given points \(\{x_i\}_{i=0}^n\) . At very point \(x_{k_j}\) of the subset, we define the following value

\[f_{x_{k_j}''}'' = \frac{2\Big[(x_{k_j} - x_{k_j-1})f(x_{k_j+1}) - (x_{k_j+1} - x_{k_j-1})f(x_{k_j}) + (x_{k_j+1} - x_{k_j})f(x_{k_j-1})\Big]}{(x_{k_j} - x_{k_j-1})(x_{k_j+1} - x_{k_j})(x_{k_j+1} - x_{k_j-1})}\,.\]

Based on the smaller set of data

\[\{(x_{k_j}, f_{x_{k_j}}'')\}_{j=1}^N, \quad h_2 := \max_{2 \le j \le N} (x_{k_j} - x_{k_{j-1}}), \tag{13}\]

by using the IMQ-RBF given in (7), a RBF interpolant \(S_{f'',N}\) is obtained, and satisfying

\[S_{f'',N}(x_{k_j}) = f''_{x_{k_i}}, \quad j = 1, \ldots, N.\]

From the above subsection, the interpolant \(S_{f'',N}\) can be expressed as follows

\[S_{f''}(x) = \sum_{i=1}^{N} \alpha_{i} \varphi(\|x - x_{k_{i}}\|), \tag{14}\]

and the coefficients \(\{\alpha_j\}_{j=1}^N\) are uniquely determined by the interpolation condition

\[S_{f''}(\mathbf{x}_{k_i}) = \sum_{j=1}^{N} \alpha_j \varphi(\|\mathbf{x}_{k_i} - \mathbf{x}_{k_j}\|) = f''(\mathbf{x}_{k_i}), \quad i = 1, \dots, N.\tag{15}\]

By using f(x) and the coefficients \(\{\alpha_i\}_{i=1}^N\) mentioned above, a function is constructed in the form

\[e(x) = f(x) - \sum_{j=1}^{N} \alpha_j \sqrt{s^2 + (x - x_{k_j})^2}.\tag{16}\]

Then a MQ quasi-interpolation can be obtained by using
\(\mathcal{L}_{\mathcal{D}}\) defined by (11) and (12) on the data \(\{x_i, e(x_i)\}_{i=0}^n\) , with the shape parameter c.~Furthermore, the MQ quasi-interpolation operator \(\mathcal{L}_{\mathcal{W}_2}\) is presented as

\[\mathcal{L}_{w_2} f(x) = \sum_{j=1}^{N} \alpha_j \sqrt{s^2 + (x - x_{k_j})^2} + \mathcal{L}_{\mathcal{D}} e(x).\tag{17}\]

Generally speaking, the shape parameters c and s should not be the same constant in (17). Moreover, the linear reproducing property and the high convergence rate of \(\mathcal{L}_{w_2}\) were also studied in {[}17{]}.
Apparently, in application one can use the second derivative of \((\mathcal{L}_{w_2}f)(x)\)

\[(\mathcal{L}_{\mathcal{W}_2} f)''(x) = \sum_{j=1}^N \alpha_j \varphi(x - x_{k_j}) + \sum_{i=0}^n \left( f(x_i) - \sum_{j=1}^N \alpha_j \sqrt{s^2 + (x_i - x_{k_j})^2} \right) \psi_i''(x),\tag{17A}\]
where \(x \in [x_0, x_n]\), to approximate \(f''(x)\).

\hypertarget{numerical-scheme-using-mq-quasi-interpolation-mathcall_w_2}{%
\paragraph[3. Numerical scheme using MQ quasi-interpolation
\(\mathcal{L}_{w_2}\)]{\texorpdfstring{\protect\hypertarget{page-3-0}{}{}3.
Numerical scheme using MQ quasi-interpolation \(\mathcal{L}_{w_2}\)}{3. Numerical scheme using MQ quasi-interpolation \textbackslash mathcal\{L\}\_\{w\_2\}}}\label{numerical-scheme-using-mq-quasi-interpolation-mathcall_w_2}}

Now, we present the numerical scheme for solving the SG equation by using the MQ quasi-interpolation \(\mathcal{L}_{w_2}\) . In our method, we use the derivative of the MQ quasi-interpolation to approximate the spatial derivative of the differential equations and employ a first
order an accurate forward difference to approximate the temporal derivative.

Discretizing the SG equation

\[\frac{\partial^2 u}{\partial t^2} = \frac{\partial^2 u}{\partial x^2} - \sin(u)\]

in time, we get

\[u_k^{n+1} = 2u_k^n - u_k^{n-1} + \tau^2(u_{xx})_k^n - \tau^2 \sin(u_k^n),\]

where \(u_k^n\) is the approximation of the value u(x,t) at
\((x_k,t_n)\) , \(t_n=n\tau\) , and \(\tau\) is the time step. Because equation (1) is a developmental equation, it is reasonable to employ the finite difference method in the time domain. Then, we use the second derivative of the MQ quasi-interpolation \(\mathcal{L}_{\mathcal{W}}\),
\(u(x,t_n)\) to approximate \(u_{xx}\) . Thus, the resulting numerical scheme is

\[u_{k}^{n+1} = 2u_{k}^{n} - u_{k}^{n-1} - \tau^{2} \sin(u_{k}^{n}) + \tau^{2} \left[ \sum_{j=1}^{N} \alpha_{j} \varphi(x_{k} - x_{k_{j}}) + \sum_{i=0}^{n} \left( u_{i}^{n} - \sum_{j=1}^{N} \alpha_{j} \sqrt{s^{2} + (x_{i} - x_{k_{j}})^{2}} \right) \psi_{i}''(x_{k}) \right], \tag{18}\]

where \(\varphi\) and \(\psi\) are defined in (7) and (12) respectively, and \(\{\alpha_j\}_{j=1}^N\) is determined by the uniquely solvable linear system

\[\sum_{j=1}^{N} \alpha_{j} \varphi(\|x_{k_{i}} - x_{k_{j}}\|) = \frac{2u_{k_{i}+1}^{n}}{(x_{k_{i}+1} - x_{k_{i}})(x_{k_{i}+1} - x_{k_{i}-1})} + \frac{2u_{k_{i}-1}^{n}}{(x_{k_{i}} - x_{k_{i}-1})(x_{k_{i}+1} - x_{k_{i}-1})} - \frac{2u_{k_{i}}^{n}}{(x_{k_{i}} - x_{k_{i}-1})(x_{k_{i}+1} - x_{k_{i}})}, \quad 1 \leqslant i \leqslant N.\tag{19}\]


We iterate this scheme to achieve the numerical solution of (1).

\hypertarget{numerical-example}{%
\paragraph{4. Numerical example}\label{numerical-example}}

In this section, we test the method by two classical examples. For the sake of simplification, we set

\[x_i = a + \frac{i(b-a)}{n}, \quad i = 0, \ldots, n,\]

and \(N = \frac{n}{K}\) , where K is a positive integer. The choice of
the interpolation centers \(\{x_{k_j}\}_{j=1}^N\) as follows

\[x_{k_1} = x_1, \quad x_{k_N} = x_{n-1},\]

\[x_{k_j} = a + \frac{(j-1)(b-a)}{N}, \quad j = 2, \dots, N.\]

The accuracy of the proposed method is measured using the
root-mean-square error (RMS) and \(L_{\infty}\) error norms defined as

\[RMS = \frac{1}{n+1} \sqrt{\sum_{i=0}^{n} |u_i^{\text{exact}} - u_i^{\text{num}}|^2},\]

and

\[L_{\infty} = \max_{0 \le i \le n} |u_i^{\text{exact}} - u_i^{\text{num}}|.\]

\textbf{Example 1.} In this example, we consider SG equation (1) without nonlinear term \(\sin(u)\) in the region \(-1 \le x \le 1\) . The initial conditions are given by

\[u(x, 0) = f(x) = \sin(\pi x),\]

\(u_t(x, 0) = g(x) = 0,\)

with the boundary conditions

\[u(-1,t) = u(1,t) = 0.\]

The analytical solution is given in {[}17{]} as

\[u(x,t) = \frac{1}{2}(\sin(\pi(x+t)) + \sin(\pi(x-t))).\]

The \(L_{\infty}\) errors and RMS errors are obtained in Table 1 for
\(t=0.25,\ 0.5,\ 0.75\) and 1. The graph of the error function
\(\widetilde{u}(x,t)-u(x,t)\) for t=1 is given in Fig. 1, where
\(\widetilde{u}(x,t)\) is the estimated function and u(x,t) is the
analytical solution.

\protect\hypertarget{page-4-0}{}{}\textbf{Table 1} Computational results
for Example 1.

\begin{longtable}[]{@{}llll@{}}
\toprule\noalign{}
t & \(L_{\infty}\) -error & RMS & Time (s) \\
\midrule\noalign{}
\endhead
\bottomrule\noalign{}
\endlastfoot
0.25 & \(1.24\times10^{-2}\) & \(6.46\times10^{-4}\) & 0.61 \\
0.50 & \(5.47\times10^{-2}\) & \(2.71\times10^{-3}\) & 0.64 \\
0.75 & \(6.14\times10^{-2}\) & \(3.04\times10^{-3}\) & 0.66 \\
1.00 & \(7.59\times10^{-3}\) & \(3.82\times10^{-4}\) & 0.67 \\
\end{longtable}

\(L_{\infty}\) and RMS errors, with
\(c=0.027,\ s=0.8,\ \tau=h=0.01,\ K=8.\)

\includegraphics{images/_page_4_Figure_5.jpeg}

\textbf{Fig. 1.} Error function \(\widetilde{u}(x,t) - u(x,t)\) in t = 1, with dt = dx = 0.01 and K = 8, for Example 1.

\textbf{Table 2} Comparison of \(L_\infty\) and RMS errors between
\(\mathcal{L}_{\mathcal{E}}\) and \(\mathcal{L}_{\mathcal{W}_2}\)
schemes.

\begin{longtable}[]{@{}
  >{\raggedright\arraybackslash}p{(\columnwidth - 8\tabcolsep) * \real{0.0388}}
  >{\raggedright\arraybackslash}p{(\columnwidth - 8\tabcolsep) * \real{0.2248}}
  >{\raggedright\arraybackslash}p{(\columnwidth - 8\tabcolsep) * \real{0.2558}}
  >{\raggedright\arraybackslash}p{(\columnwidth - 8\tabcolsep) * \real{0.2248}}
  >{\raggedright\arraybackslash}p{(\columnwidth - 8\tabcolsep) * \real{0.2558}}@{}}
\toprule\noalign{}
\begin{minipage}[b]{\linewidth}\raggedright
t
\end{minipage} & \begin{minipage}[b]{\linewidth}\raggedright
\(L_{\infty}\) error
\end{minipage} & \begin{minipage}[b]{\linewidth}\raggedright
\end{minipage} & \begin{minipage}[b]{\linewidth}\raggedright
RMS
\end{minipage} & \begin{minipage}[b]{\linewidth}\raggedright
\end{minipage} \\
\midrule\noalign{}
\endhead
\bottomrule\noalign{}
\endlastfoot
& \(\mathcal{L}_{\mathcal{E}}\) & \(\mathcal{L}_{\mathcal{W}_2}^*\) &
\(\mathcal{L}_{\mathcal{E}}\) & \(\mathcal{L}_{\mathcal{W}_2}^*\) \\
0.1 & \(1.54 \times 10^{-6}\) & \(1.31 \times 10^{-5}\) &
\(7.43 \times 10^{-6}\) & \(3.39\times10^{-7}\) \\
0.2 & \(4.25\times10^{-5}\) & \(2.50\times10^{-5}\) &
\(1.76\times10^{-5}\) & \(6.55\times10^{-7}\) \\
0.3 & \(9.02\times10^{-5}\) & \(3.47\times10^{-5}\) &
\(3.60\times10^{-5}\) & \(9.29\times10^{-7}\) \\
0.4 & \(1.62\times10^{-4}\) & \(4.20\times10^{-5}\) &
\(1.62\times10^{-4}\) & \(1.15\times10^{-6}\) \\
0.5 & \(2.58\times10^{-4}\) & \(4.68 \times 10^{-5}\) &
\(1.10 \times 10^{-4}\) & \(1.34\times10^{-6}\) \\
0.6 & \(3.73 \times 10^{-4}\) & \(6.54\times10^{-5}\) &
\(1.65\times10^{-4}\) & \(1.55\times10^{-6}\) \\
0.7 & \(4.98\times10^{-4}\) & \(1.74\times10^{-4}\) &
\(2.29\times10^{-4}\) & \(2.14\times10^{-6}\) \\
0.8 & \(6.24\times10^{-4}\) & \(4.01\times10^{-4}\) &
\(2.98\times10^{-4}\) & \(3.92\times10^{-6}\) \\
0.9 & \(7.44\times10^{-4}\) & \(8.22\times10^{-4}\) &
\(3.69 \times 10^{-4}\) & \(7.98\times10^{-6}\) \\
1.0 & \(8.49\times10^{-4}\) & \(1.53 \times 10^{-3}\) &
\(4.37\times10^{-4}\) & \(1.56\times10^{-5}\) \\
\end{longtable}

* denotes that we set K = 10, s = 1.

\textbf{Example 2.} In this example, we compare the accuracy of the MQ quasi-interpolation \(\mathcal{L}_{w_2}\) proposed in this paper to the scheme \(\mathcal{L}_{\mathcal{E}}\) used by Ma and Wu, thus we utilize the same example adopted in {[}13{]}.

We consider SG equation (1) in the domain \(-2 \le x \le 2\) , and the initial conditions are given by

\[u(x, 0) = g(x) = 0,\]

\(u_t(x, 0) = h(x) = 4\operatorname{sech}(x).\)

The analytical solution is

\[u(x, t) = 4 \arctan(\operatorname{sech}(x)t).\]

For comparison, the numerical results are commutated with parameters:
time step \(\tau=0.01\) , space step h=0.01 and shape parameter \(c=0.1\sqrt[3]{h}\) . Table 2 show the values of \(L_{\infty}\) and RMS errors obtained by using \(\mathcal{L}_{\mathcal{E}}\) and
\(\mathcal{L}_{\mathcal{W}_2}\) quasi-interpolation schemes, at \(t=0.1,\ 0.2,\ \ldots,\ 1.00\) respectively.

\protect\hypertarget{page-5-0}{}{}From Table 2, we see that the solutions found with \(\mathcal{L}_{w_2}\) scheme is in good agreement with the results obtained by \(\mathcal{L}_{\mathcal{E}}\) scheme in the sense of \(I_{\infty}\) errors, and better than it in the sense of RMS
errors.

\hypertarget{conclusions}{%
\paragraph{5. Conclusions}\label{conclusions}}

In this paper, a numerical scheme for the one-dimensional nonlinear Sine--Gordon equation by high accuracy MQ quasiinterpolation \(\mathcal{L}_{\mathcal{W}_2}\) is presented. Form the examples, we can say that the \(\mathcal{L}_{\mathcal{W}_2}\) scheme is feasible and the
accuracy is better than Ma and Wu's \(\mathcal{L}_{\mathcal{E}}\) scheme in the sense of RMS errors.

For comparison, our technique is used for the non-equidistant grids though we use equidistant grids in our numerical experiments. We see that the present technique requires calculating \(\psi_i(x)\) , \(\psi n_i(x)\) \((i=0,\ldots,n)\) , and the inverse matrix of
\((\varphi(\|x_{k_i}-x_{k_i}\|))_{N\times N}\) only once.

\hypertarget{references}{%
\paragraph{References}\label{references}}

\begin{itemize}
\tightlist
\item
  {[}1{]} J. Argyris, M. Haase, J.C. Heinrich, Finite element
  approximation to two-dimensional Sine--Gordon solutions, Comput.
  Methods Appl. Mech. Eng. 86 (1991) 1--26.
\item
  {[}2{]} R.K. Beatson, M.J.D. Powell, Univariate multiquadric
  approximation: quasi-interpolation to scattered data, Constr. Approx.
  8 (1992) 275-288.
\item
  {[}3{]} A.G. Bratsos, E.H. Twizell, A family of parametric
  finite-difference methods for the solution of the Sine-Gordon
  equation, Appl. Math. Comput. 93 (1998) 117--137.
\item
  {[}4{]} A.G. Bratsos, An explicit numerical scheme for the Sine-Gordon  equation in 2+1 dimensions, Appl. Numer. Anal. Comput. Math. 2  (2005) 189-211.
\item
  {[}5{]} M.D.~Buhmann, Multivariate interpolation with radial basis  functions, Report DAMTP 1988/NA8, University of Cambridge.
\item
  {[}6{]} M.D.~Buhmann, Radial Basis Functions: Theory and
  Implementations, Cambridge University Press, 2003.
\item
  {[}7{]} M. Dehghan, A. Shokri, A numerical method for one-dimensional
  nonlinear Sine--Gordon equation using collocation and radial basis
  functions, Numer. Methods Part D. E. 24 (2007) 687--698.
\item
  {[}8{]} B.Y. Guo, P.J. Pascual, M.J. Rodriguez, L. Vàzquez, Numerical
  solution of the Sine-Gordon equation, Appl. Math. Comput. 18 (1986)
  1-14.
\item
  {[}9{]} Z.W. Jiang, R.H. Wang, C.G. Zhu, M. Xu, High accuracy
  multiquadric quasi-interpolation, Appl. Math. Model. 35 (2011)
  2185-2195.
\item
  {[}10{]} D.G. Krige, A statistical approach to some mine valuation and
  allied problems on the Witwatersrand, M.Sc. Thesis, University of
  Witwatersrand, South Africa, 1951.
\item
  {[}11{]} J. Kuang, L. Lu, Two classes of finite-difference methods for
  generalized Sine-Gordon equations, J. Comput. Appl. Math. 31 (1990)
  389-396.
\item
  {[}12{]} L. Ling, A univariate quasi-multiquadric interpolation with
  better smoothness, Comput. Math. Appl. 48 (2004) 897-912.
\item
  {[}13{]} M.L. Ma, Z.M. Wu, A numerical method for one-dimensional
  nonlinear Sine-Gordon equation using multiquadric quasi-interpolation,
  Chin. Phys. B 18 (2009) 3099--3103.
\item
  {[}14{]} C.A. Micchelli, Interpolation of scattered data: distance
  matrices and conditionally positive definite functions, Constr.
  Approx. 2 (1986) 11--22.
\item
  {[}15{]} J.K. Perring, T.H. Skyrme, A model unified field equation,
  Nucl. Phys. 31 (1962) 550-555.
\item
  {[}16{]} J. Rashidinia, M. Ghasemia, R. Jalilian, Numerical solution
  of the nonlinear Klein-Gordon equation, J. Comput. Appl. Math. 233
  (2010) 1866-1878.
\item
  {[}17{]} Y. Wang, B. Wang, High-order multi-symplectic schemes for the
  nonlinear Klein-Gordon equation, Appl. Math. Comput. 166 (2005)
  608-632.
\item
  {[}18{]} H. Wendland, Scattered Data Approximation, Cambridge
  Monographs on Applied and Computational Mathematics, vol.~17,
  Cambridge University Press, Cambridge, 2005.
\item
  {[}19{]} G.B. Whitham, Linear and Nonlinear Waves, Wiley-Interscience,
  NewYork, 1999.
\item
  {[}20{]} Z.M. Wu, R. Schaback, Shape preserving properties and
  convergence of univariate multiquadric quasi-interpolation, Acta Math.
  Appl. Sin. 10 (4) (1994) 441--446.
\end{itemize}

\end{document}
